\documentclass[12pt]{article}
\usepackage[margin=0.75in]{geometry}
\usepackage{hyperref}
\usepackage{fontspec}
\usepackage{amsmath}
\usepackage{amssymb}
\usepackage{enumitem}
\usepackage{listings}
\usepackage{xcolor}
\usepackage{booktabs}
\usepackage{array}
\usepackage{multicol}

% Color scheme
\definecolor{trekblue}{HTML}{1A73E8}
\definecolor{trekadmon}{HTML}{D0021F}
\definecolor{trekgray}{HTML}{F5F5F5}

\hypersetup{
    colorlinks=true,
    linkcolor=trekblue,
    filecolor=trekblue,
    urlcolor=trekblue,
}

\setlist[itemize]{leftmargin=*,nosep,label=$\bullet$}
\setlist[enumerate]{leftmargin=*,nosep}

\title{\textbf{Trek Agent Cheatsheet}}
\author{@trek/coding-agent}
\date{\today}

\begin{document}

\maketitle

\section*{Overview}

The \texttt{@trek/coding-agent} is a terminal-based coding agent that extends Pi with custom tools, skills, and extensions. This cheatsheet summarizes the most useful commands, keybindings, and features.

\section{Quick Start}

\begin{itemize}
    \item Install: \texttt{npm install -g --ignore-scripts @trek/coding-agent}
    \item Authenticate: \texttt{export ANTHROPIC\_API\_KEY=sk-ant-...} then \texttt{pi}
    \item Or login via subscription: \texttt{/login}
    \item Run with initial prompt: \texttt{pi "List all .ts files in src/"}
    \item Non-interactive print mode: \texttt{pi -p "Summarize this codebase"}
\end{itemize}

\section{Interactive Commands}

Type \texttt{/} in the editor to trigger commands.

\begin{multicols}{2}
\begin{tabular}{@{}ll@{}}
\toprule
\textbf{Command} & \textbf{Description} \\
\midrule
\texttt{/login}, \texttt{/logout} & OAuth authentication \\
\texttt{/model} & Switch models \\
\texttt{/scoped-models} & Enable/disable models for Ctrl+P cycling \\
\texttt{/settings} & Thinking level, theme, message delivery, transport \\
\texttt{/thinking briefly|deep|off} & Pin short-think, restore deep thinking, or turn thinking off \\
\texttt{/hierarchy} & Layering, mode, suite, and one role \\
\texttt{/resume} & Pick from previous sessions \\
\texttt{/new} & Start a new session \\
\texttt{/name <name>} & Set session display name \\
\texttt{/session} & Show session info (file, ID, messages, tokens, cost) \\
\texttt{/tree} & Jump to any point in the session and continue from there \\
\texttt{/fork} & Create a new session from a previous user message \\
\texttt{/clone} & Duplicate the current active branch into a new session \\
\texttt{/compact [prompt]} & Manually compact context, optional custom instructions \\
\texttt{/copy} & Copy last assistant message to clipboard \\
\texttt{/export [file]} & Export session to HTML file \\
\texttt{/share} & Upload as private GitHub gist with shareable HTML link \\
\texttt{/reload} & Reload keybindings, extensions, skills, prompts, and context files \\
\texttt{/hotkeys} & Show all keyboard shortcuts \\
\texttt{/changelog} & Display version history \\
\texttt{/quit} & Quit pi \\
\bottomrule
\end{tabular}
\end{multicols}

\section{Keyboard Shortcuts}

\subsection*{Editor}
\begin{tabular}{@{}ll@{}}
\toprule
\textbf{Key} & \textbf{Action} \\
\midrule
Ctrl+C & Clear editor (twice to quit) \\
Escape & Cancel / abort \\
Escape twice & Open \texttt{/tree} \\
Ctrl+L & Open model selector \\
Ctrl+P / Shift+Ctrl+P & Cycle scoped models forward/backward \\
Shift+Tab & Cycle thinking level \\
Ctrl+O & Collapse/expand tool output \\
Ctrl+T & Collapse/expand thinking blocks \\
Shift+Enter & Insert new line (multi-line) \\
Ctrl+V & Paste image from clipboard (Alt+V on Windows) \\
\texttt{!command} & Run bash command, send output to LLM \\
\texttt{!!command} & Run bash command without sending output \\
@file & Fuzzy-search project files \\
\bottomrule
\end{tabular}

\subsection*{Sessions \& Models}
\begin{tabular}{@{}ll@{}}
\toprule
\textbf{Key} & \textbf{Action} \\
\midrule
Ctrl+L & Open model selector \\
Ctrl+P & Cycle to next model \\
Alt+P & Cycle to previous model \\
Shift+Tab & Cycle thinking level \\
Ctrl+T & Collapse/expand thinking blocks \\
Ctrl+O & Expand/collapse tool/thinking cards \\
Alt+Enter & Queue follow-up message \\
Alt+Up & Restore queued messages to editor \\
\bottomrule
\end{tabular}

\section{Built-in Tools}

\begin{tabular}{@{}ll@{}}
\toprule
\textbf{Tool} & \textbf{Description} \\
\midrule
\texttt{read} & Read file contents \\
\texttt{write} & Create or overwrite files \\
\texttt{edit} & Precise text replacement in files \\
\texttt{bash} & Execute shell commands \\
\texttt{grep} & Search for patterns in files \\
\texttt{find} & Find files by name or pattern \\
\texttt{ls} & List directory contents \\
\bottomrule
\end{tabular}

\section{Session Management}

\begin{tabular}{@{}ll@{}}
\toprule
\textbf{CLI Flag} & \textbf{Description} \\
\midrule
\texttt{pi -c} & Continue most recent session \\
\texttt{pi -r} & Browse and select from past sessions \\
\texttt{pi --no-session} & Ephemeral mode (don't save) \\
\texttt{pi --name "task"} & Set session display name at startup \\
\texttt{pi --session <path|id>} & Use specific session file or ID \\
\texttt{pi --fork <path|id>} & Fork session into a new session \\
\bottomrule
\end{tabular}

\section{Settings}

Use \texttt{/settings} to modify common options, or edit JSON files directly:

\begin{itemize}
    \item Global: \texttt{~/.trek/agent/settings.json}
    \item Project: \texttt{.trek/settings.json} (overrides global)
\end{itemize}

\subsection*{Environment Variables}
\begin{tabular}{@{}ll@{}}
\toprule
\textbf{Variable} & \textbf{Description} \\
\midrule
\texttt{PI\_OFFLINE} & Disable startup network operations \\
\texttt{PI\_SKIP\_VERSION\_CHECK} & Skip version update check at startup \\
\texttt{PI\_TELEMETRY} & Override install/update telemetry (0/1) \\
\texttt{PI\_CACHE\_RETENTION} & Set to \texttt{long} for extended prompt cache \\
\texttt{VISUAL}, \texttt{EDITOR} & External editor for Ctrl+G \\
\bottomrule
\end{tabular}

\section{Customization}

\subsection*{Prompt Templates}
Reusable prompts as Markdown files. Type \texttt{/name} to expand.
\begin{itemize}
    \item Place in \texttt{~/.trek/agent/prompts/}, \texttt{.pi/prompts/}, or a pi package
    \item Example: \texttt{Review this code for bugs. Focus on: \{ \{focus\} \}}
\end{itemize}

\subsection*{Skills}
On-demand capability packages. Invoke via \texttt{/skill:name}.
\begin{itemize}
    \item Place in \texttt{~/.trek/agent/skills/}, \texttt{.pi/skills/}, or a pi package
    \item See \texttt{docs/skills.md}
\end{itemize}

\subsection*{Extensions}
TypeScript modules that extend pi with custom tools, commands, keyboard shortcuts, event handlers, and UI components.
\begin{itemize}
    \item Place in \texttt{~/.trek/agent/extensions/}, \texttt{.pi/extensions/}, or a pi package
    \item See \texttt{docs/extensions.md}
\end{itemize}

\subsection*{Themes}
Built-in: \texttt{dark}, \texttt{light}. Themes hot-reload automatically.
\begin{itemize}
    \item Place in \texttt{~/.trek/agent/themes/}, \texttt{.pi/themes/}, or a pi package
    \item See \texttt{docs/themes.md}
\end{itemize}

\subsection*{Pi Packages}
Bundle and share extensions, skills, prompts, and themes via npm or git.
\begin{itemize}
    \item \texttt{pi install npm:@foo/pi-tools}
    \item \texttt{pi install git:github.com/user/repo@v1}
    \item \texttt{pi list} \quad \texttt{pi update} \quad \texttt{pi remove <source>}
\end{itemize}

\section{CLI Reference}

\begin{tabular}{@{}ll@{}}
\toprule
\textbf{Flag} & \textbf{Description} \\
\midrule
\texttt{-p}, \texttt{--print} & Print response and exit \\
\texttt{--mode json} & Output all events as JSON lines \\
\texttt{--mode rpc} & RPC mode for process integration \\
\texttt{--export <in> [out]} & Export session to HTML \\
\texttt{--provider <name>} & Provider (anthropic, openai, google, etc.) \\
\texttt{--model <pattern>} & Model pattern or ID \\
\texttt{--api-key <key>} & API key (overrides env vars) \\
\texttt{--thinking <level>} & \texttt{off}, \texttt{minimal}, \texttt{low}, \texttt{medium}, \texttt{high}, \texttt{xhigh} \\
\texttt{--models <patterns>} & Comma-separated patterns for Ctrl+P cycling \\
\texttt{--list-models [search]} & List available models \\
\texttt{-t}, \texttt{--tools <list>} & Allowlist specific tool names \\
\texttt{-xt}, \texttt{--exclude-tools <list>} & Disable specific tool names \\
\texttt{-nbt}, \texttt{--no-builtin-tools} & Disable built-in tools by default \\
\texttt{-nt}, \texttt{--no-tools} & Disable all tools by default \\
\texttt{-e}, \texttt{--extension <source>} & Load extension from path, npm, or git \\
\texttt{--no-extensions} & Disable extension discovery \\
\texttt{--skill <path>} & Load skill (repeatable) \\
\texttt{--no-skills} & Disable skill discovery \\
\texttt{--prompt-template <path>} & Load prompt template (repeatable) \\
\texttt{--no-prompt-templates} & Disable prompt template discovery \\
\texttt{--theme <path>} & Load theme (repeatable) \\
\texttt{--no-themes} & Disable theme discovery \\
\texttt{-nc}, \texttt{--no-context-files} & Disable AGENTS.md and CLAUDE.md discovery \\
\texttt{--system-prompt <text>} & Replace default prompt \\
\texttt{--append-system-prompt <text>} & Append to system prompt \\
\texttt{--verbose} & Force verbose startup \\
\texttt{-h}, \texttt{--help} & Show help \\
\texttt{-v}, \texttt{--version} & Show version \\
\bottomrule
\end{tabular}

\section{File Arguments}

Prefix files with \texttt{@} to include in the message:
\begin{itemize}
    \item \texttt{pi @prompt.md "Answer this"}
    \item \texttt{pi -p @screenshot.png "What's in this image?"}
    \item \texttt{pi @code.ts @test.ts "Review these files"}
\end{itemize}

\section{Philosophy}

\begin{itemize}
    \item \textbf{No MCP.} Build CLI tools with READMEs, or build an extension that adds MCP support.
    \item \textbf{No sub-agents.} Spawn pi instances via tmux, or build your own with extensions.
    \item \textbf{No permission popups.} Run in a container, or build your own confirmation flow with extensions.
    \item \textbf{No plan mode.} Write plans to files, or build it with extensions.
    \item \textbf{No built-in to-dos.} Use a TODO.md file, or build your own with extensions.
\end{itemize}

\end{document}
